% ============================================================
%  APUNTES DE LECTOCOMPRENSIÓN — CLÁUSULAS ESTÁNDAR EN CONTRATOS
%  Autor: Isaac Edgar Camacho Ocampo
%  Septiembre 2026
% ============================================================

\documentclass[12pt,oneside]{book}

% ────────────────────────────────────────────────────────────
%  METADATOS
% ────────────────────────────────────────────────────────────
\newcommand{\universidad}{Universidad de Buenos Aires (UBA)}
\newcommand{\facultad}{}
\newcommand{\carrera}{}
\newcommand{\materia}{Lectocomprensión}
\newcommand{\titulo}{Cláusulas Estándar en Contratos (Boilerplate Clauses)}
\newcommand{\autor}{Isaac Edgar Camacho Ocampo}
\newcommand{\anio}{2026}

% ────────────────────────────────────────────────────────────
%  PAQUETES
% ────────────────────────────────────────────────────────────
\usepackage[utf8]{inputenc}
\usepackage[T1]{fontenc}
\usepackage[spanish]{babel}

\usepackage[
    top=2.5cm,
    bottom=2.5cm,
    left=3cm,
    right=2.5cm,
    headheight=15pt
]{geometry}

\usepackage{amsmath}
\usepackage{amssymb}
\usepackage{mathtools}

\usepackage{graphicx}
\usepackage{float}
\usepackage{caption}
\usepackage{subcaption}

\usepackage{booktabs}
\usepackage{array}
\usepackage{longtable}
\usepackage{tabularx}

\usepackage{enumitem}

\usepackage{xcolor}
\usepackage[most]{tcolorbox}

\usepackage{fancyhdr}
\usepackage{ragged2e}

\usepackage[
    colorlinks=true,
    linkcolor=blue!50!black,
    citecolor=green!40!black,
    urlcolor=blue!60!black,
    pdftitle={Apuntes de Lectocomprensión — Cláusulas Estándar en Contratos},
    pdfauthor={Isaac Edgar Camacho Ocampo},
    pdfsubject={Apuntes universitarios},
    bookmarks=true
]{hyperref}

% ────────────────────────────────────────────────────────────
%  RUTAS DE IMÁGENES
% ────────────────────────────────────────────────────────────
\graphicspath{
    {./img/}
    {./graficos/}
    {./figuras/}
}

% ────────────────────────────────────────────────────────────
%  CONFIGURACIÓN DE ÍNDICE
% ────────────────────────────────────────────────────────────
\setcounter{tocdepth}{2}

% ────────────────────────────────────────────────────────────
%  CONFIGURACIÓN DE PÁRRAFOS
% ────────────────────────────────────────────────────────────
\setlength{\parindent}{0pt}
\setlength{\parskip}{0.5em}

% ────────────────────────────────────────────────────────────
%  CONFIGURACIÓN DE LISTAS
% ────────────────────────────────────────────────────────────
\setlist[itemize]{topsep=0.3em, itemsep=0.2em}
\setlist[enumerate]{topsep=0.3em, itemsep=0.2em}

% ────────────────────────────────────────────────────────────
%  CONTROL DE VIUDAS Y HUÉRFANAS
% ────────────────────────────────────────────────────────────
\clubpenalty=10000
\widowpenalty=10000

% ────────────────────────────────────────────────────────────
%  ENCABEZADOS Y PIES
% ────────────────────────────────────────────────────────────
\pagestyle{fancy}
\fancyhf{}
\fancyhead[L]{\nouppercase{\leftmark}}
\fancyhead[R]{\materia}
\fancyfoot[C]{\thepage}
\renewcommand{\headrulewidth}{0.4pt}
\renewcommand{\footrulewidth}{0pt}

\fancypagestyle{plain}{
    \fancyhf{}
    \fancyfoot[C]{\thepage}
    \renewcommand{\headrulewidth}{0pt}
}

% ────────────────────────────────────────────────────────────
%  COMANDOS MATEMÁTICOS
% ────────────────────────────────────────────────────────────
\newcommand{\abs}[1]{\left|#1\right|}
\newcommand{\norm}[1]{\left\lVert#1\right\rVert}
\newcommand{\vect}[1]{\mathbf{#1}}
\newcommand{\deriv}[2]{\frac{d#1}{d#2}}
\newcommand{\pderiv}[2]{\frac{\partial#1}{\partial#2}}

% ────────────────────────────────────────────────────────────
%  CAJAS ACADÉMICAS
% ────────────────────────────────────────────────────────────

% Definición
\newtcolorbox{definition}[1]{
    enhanced,
    breakable,
    title={Definición: #1},
    fonttitle=\bfseries,
    colback=blue!3,
    colframe=blue!50!black,
    boxrule=0.8pt,
    arc=2mm
}

% Importante
\newtcolorbox{important}{
    enhanced,
    breakable,
    title={Importante},
    fonttitle=\bfseries,
    colback=yellow!8,
    colframe=orange!70!black,
    boxrule=0.8pt,
    arc=2mm
}

% Ejemplo
\newtcolorbox{example}[1]{
    enhanced,
    breakable,
    title={Ejemplo: #1},
    fonttitle=\bfseries,
    colback=green!4,
    colframe=green!40!black,
    boxrule=0.8pt,
    arc=2mm
}

% Fórmula / Regla
\newtcolorbox{formula}[1]{
    enhanced,
    breakable,
    title={Fórmula / Regla: #1},
    fonttitle=\bfseries,
    colback=purple!4,
    colframe=purple!50!black,
    boxrule=0.8pt,
    arc=2mm
}

% Ejercicio
\newtcolorbox{exercise}[1]{
    enhanced,
    breakable,
    title={Ejercicio: #1},
    fonttitle=\bfseries,
    colback=gray!5,
    colframe=gray!60!black,
    boxrule=0.8pt,
    arc=2mm
}

% Nota
\newtcolorbox{note}{
    enhanced,
    breakable,
    title={Nota},
    fonttitle=\bfseries,
    colback=white,
    colframe=gray!50,
    boxrule=0.6pt,
    arc=2mm
}

% ============================================================
%  DOCUMENTO
% ============================================================
\begin{document}

% ────────────────────────────────────────────────────────────
%  PORTADA
% ────────────────────────────────────────────────────────────
\begin{titlepage}
\thispagestyle{empty}

\begin{center}

\vspace*{1cm}

{\large \textsc{\universidad}}

\vspace{3cm}

{\Huge \textbf{\materia}}

\vspace{1cm}

{\LARGE \titulo}

\vspace{0.8cm}

\textit{Comprender el lenguaje jurídico es la primera herramienta de todo profesional del derecho.}

\vspace{1.2cm}

\rule{0.70\textwidth}{0.5pt}

\vspace{0.5cm}

{\large Apuntes de estudio}

\vspace{0.5cm}

\rule{0.70\textwidth}{0.5pt}

\vfill

{\large \autor}

\vspace{0.3cm}

{\large \anio}

\end{center}

\vfill

\begin{flushleft}
\textit{Este es un modesto aporte a los estudiantes de ``Lectocomprensión'', no pretende sustituir el material oficial de la materia.}
\end{flushleft}

\end{titlepage}

% ────────────────────────────────────────────────────────────
%  ÍNDICE
% ────────────────────────────────────────────────────────────
\tableofcontents
\clearpage

% ============================================================
%  CAPÍTULO 1 — CONCEPTOS FUNDAMENTALES
% ============================================================
\chapter{Conceptos Fundamentales: Las Cláusulas Estándar}

\section{¿Qué son las Boilerplate Clauses?}

Las \textbf{boilerplate clauses} (cláusulas estándar o cláusulas tipo) son
disposiciones contractuales estandarizadas que se replican con escasas
variaciones de un contrato a otro. Debido a su carácter genérico, suelen
ser frecuentemente pasadas por alto por las partes al momento de negociar y
suscribir un acuerdo. Sin embargo, tienen una relevancia jurídica
significativa, pues regulan cuestiones vinculadas a posibles conflictos
entre las partes.

En la práctica argentina, estas cláusulas suelen agruparse bajo el título
\textbf{``Disposiciones Varias''}, que funciona como categoría residual
para todo aquello que no constituye una cláusula específica del contrato.

\subsection{Características generales}

\begin{itemize}
    \item Son lo suficientemente estándar para no ser objeto de negociación.
    \item Abordan cuestiones vinculadas a posibles conflictos entre las partes.
    \item Están redactadas en lenguaje claro, coherente y estandarizado.
    \item La redacción suele ser muy similar de un contrato a otro.
    \item Se replican de un acuerdo a otro sin demasiada variación.
    \item Aparecen con frecuencia independientemente del tipo contractual.
    \item Suelen ubicarse al final del contrato.
\end{itemize}

\section{Origen Histórico del Término ``Boilerplate''}

\begin{definition}{Boilerplate}
El término \textbf{``boilerplate''} proviene de la imprenta del siglo XIX.
Hace referencia a las placas metálicas (de acero) que se utilizaban en las
prensas tipográficas de la época. Estas placas contenían textos
preestablecidos que se replicaban idénticamente en cada impresión, creando
múltiples copias de un mismo molde. La idea de \emph{texto que no se
modifica} y \emph{texto fijo} es la que ha perdurado en el lenguaje
jurídico moderno.
\end{definition}

\begin{important}
El término \textbf{``boilerplate''} no debe traducirse literalmente como
``placa de fundición'' o ``placa caliente''. En el lenguaje jurídico ha
evolucionado para significar texto estandarizado, texto fijo, texto
frecuente y texto que no está sujeto a negociación.
\end{important}

La relevancia de esta terminología radica en que el carácter genérico del
nombre contribuye a que estas cláusulas sean \emph{pasadas por alto}. En
el texto analizado en clase (de origen australiano), se señala que
\textit{``due to their name, boilerplate clauses can often be overlooked''}
(``debido a su nombre, estas cláusulas suelen ser omitidas'').

% ============================================================
%  CAPÍTULO 2 — JURISDICCIÓN Y DERECHO APLICABLE
% ============================================================
\chapter{Jurisdicción y Derecho Aplicable}

\section{Governing Law (Derecho Aplicable)}

\begin{definition}{Governing Law / Derecho Aplicable}
La cláusula de \textbf{derecho aplicable} (también denominada
\textit{governing law}) determina con qué legislación se interpretará e
integrará el contrato. Establece los parámetros normativos dentro de los
cuales se definirán las obligaciones, los derechos y las reglas del acuerdo.
\end{definition}

Esta cláusula es especialmente relevante en contratos internacionales, donde
las partes provienen de distintas jurisdicciones o existe algún elemento
extranjero. Ante esa situación, resulta indispensable determinar cuál será
la ley aplicable.

\section{Jurisdiction (Competencia)}

La palabra inglesa \textbf{``jurisdiction''} constituye un \textbf{falso
cognado parcial} respecto del término español ``jurisdicción''. Si bien en
ocasiones puede referirse efectivamente a la jurisdicción, en la mayoría de
los contextos contractuales alude a la \textbf{competencia}.

\begin{definition}{Jurisdicción vs. Competencia}
\begin{itemize}
    \item \textbf{Jurisdicción:} Concepto abstracto. Potestad o capacidad
    que poseen los jueces en virtud de su investidura y de su cargo para
    decidir sobre una causa. Todos los jueces tienen jurisdicción.
    \item \textbf{Competencia:} Concepto práctico. Determina qué juez
    específico, en qué zona, en qué localidad o en qué lugar debe entender
    en una causa. Se clasifica según: materia, territorio, grado y persona.
\end{itemize}
\end{definition}

Una expresión que sintetiza esta relación es: \textit{``la competencia es
la medida de la jurisdicción''}: todos los jueces tienen jurisdicción,
pero es necesario determinar a cuál de ellos le corresponde conocer en
cada caso particular.

En inglés existe una única palabra (\textit{jurisdiction}) para ambos
conceptos, lo que genera dificultades de interpretación en la lectura de
textos jurídicos en ese idioma. Al encontrar la palabra
\textit{jurisdiction} en un contrato, debe evaluarse el contexto para
determinar si hace referencia a la jurisdicción en sentido abstracto o a
la competencia.

\section{Coherencia entre Derecho Aplicable y Competencia}

Las cláusulas de derecho aplicable y de competencia deben ser coherentes
entre sí (\textbf{aligned}, es decir, alineadas). No resulta válido
establecer que el contrato se rige por la legislación de un país y que los
conflictos se someterán a la competencia de los tribunales de otro, si tal
combinación carece de coherencia jurídica o funcional.

\begin{example}{Inconsistencia entre derecho aplicable y competencia}
Sería incoherente establecer que el contrato se rige por la legislación de
Estados Unidos y, al mismo tiempo, que cualquier controversia se someterá
a los tribunales de la República Argentina. Ambas cláusulas deben
articularse de manera consistente.
\end{example}

% ============================================================
%  CAPÍTULO 3 — RIESGO Y CUMPLIMIENTO
% ============================================================
\chapter{Riesgo y Cumplimiento}

\section{Entire Agreement (Acuerdo Total)}

\begin{definition}{Entire Agreement}
La cláusula de \textbf{entire agreement} (acuerdo total o acuerdo
íntegro) establece que únicamente lo consignado por escrito en el
contrato es vinculante para las partes. Las negociaciones previas,
las comunicaciones (correos electrónicos, llamadas telefónicas) y las
manifestaciones realizadas durante las tratativas anteriores a la
suscripción \emph{no} integran el acuerdo ni tienen fuerza obligatoria.
\end{definition}

El contrato en sí mismo define el marco dentro del cual se desarrollará
la relación entre las partes: sus obligaciones, sus derechos y las pautas
de funcionamiento del acuerdo.

\begin{note}
Es habitual que, previo a la firma de un contrato, hayan tenido lugar
negociaciones: llamadas telefónicas, intercambio de correos electrónicos,
reuniones. La cláusula de \textit{entire agreement} establece
expresamente que esas comunicaciones y manifestaciones previas no serán
consideradas parte del contrato.
\end{note}

\subsection{Comparación con el sistema argentino}

En el derecho argentino esta cláusula resulta en principio innecesaria,
dado que el principio de \textbf{autonomía de la voluntad} está implícito
en el sistema. El Código Civil y Comercial establece que, ante cláusulas
poco claras o controvertidas, se acudirá primero a la voluntad de las
partes como criterio interpretativo y, subsidiariamente, a las costumbres
desarrolladas durante la celebración del contrato. Lo que no consta en el
instrumento no es vinculante para las partes.

Sin embargo, por efecto de la globalización y del comercio internacional,
esta cláusula comienza a incorporarse en contratos locales, adoptada de
modelos anglosajones.

\section{Force Majeure (Fuerza Mayor)}

\begin{definition}{Force Majeure}
La \textbf{fuerza mayor} (\textit{force majeure}) designa un evento o
circunstancia que está fuera del control de las partes, que no es
responsabilidad de ninguna de ellas y que impide el cumplimiento de las
obligaciones contractuales. Funciona como una justificación o mecanismo
de anticipación ante el incumplimiento derivado de causas ajenas a la
voluntad de los contratantes.
\end{definition}

\subsection{Función de la cláusula}

La cláusula de fuerza mayor establece los \textbf{grounds} (causales)
por los cuales una parte puede quedar eximida de su obligación de cumplir.
El término jurídico \textit{grounds} significa ``causales'': por ejemplo,
\textit{grounds for divorce} equivale a ``causales de divorcio''. En el
contrato, la cláusula tiende a detallar de manera exhaustiva cuáles son
esas causales habilitantes.

\begin{example}{Evento de fuerza mayor}
La pandemia de COVID-19 fue un evento paradigmático de fuerza mayor en el
ámbito contractual. Generó situaciones de incumplimiento derivadas de
circunstancias completamente ajenas a la voluntad de las partes y que no
podían haber sido razonablemente previstas.
\end{example}

\subsection{Diferencia con el sistema argentino}

En los contratos redactados en inglés no existe, en general, una
definición legal predeterminada de fuerza mayor, por lo que la cláusula
debe ser incluida expresamente y debe definir con precisión qué
circunstancias la configuran. En el derecho argentino, en cambio, el
Código Civil y Comercial define la fuerza mayor, de modo que no es
necesaria una definición exhaustiva en el texto del contrato.

El sistema argentino cuenta además con la \textbf{Teoría de la
Imprevisión}, que contempla circunstancias sobrevinientes, no previstas ni
previsibles al momento de la celebración, que afectan el desarrollo normal
del cumplimiento de las obligaciones.

\section{Acts of God (Actos de Dios)}

\begin{definition}{Acts of God vs. Force Majeure}
\begin{itemize}
    \item \textbf{Acts of God:} Eventos de la naturaleza que ocurren sin
    intervención humana (terremotos, inundaciones, huracanes, etc.).
    \item \textbf{Force Majeure:} Concepto más amplio que abarca tanto los
    actos de la naturaleza como eventos de origen humano (guerras, disturbios
    civiles, huelgas, etc.).
\end{itemize}
Ambos conceptos son equivalentes, en términos generales, a la fuerza mayor
y al caso fortuito del derecho argentino. La \textit{force majeure} es la
categoría más cercana al caso fortuito; los \textit{acts of God} serían su
subconjunto referido exclusivamente a eventos naturales.
\end{definition}

\section{Termination (Extinción Anticipada) y Discharge (Cumplimiento)}

\begin{definition}{Terminate vs. Discharge}
\begin{itemize}
    \item \textbf{Terminate:} Poner fin al contrato de manera anticipada.
    Equivale a las formas anormales de extinción del derecho argentino
    (rescisión, resolución, revocación, nulidad, novación). La expresión
    \textit{early termination} (extinción temprana/anticipada) refuerza
    este sentido.
    \item \textbf{Discharge:} Concepto más general. Abarca tanto el
    cumplimiento o el vencimiento del plazo (extinción normal) como otras
    formas de extinción. Cuando el contrato se extinguió simplemente por
    el transcurso del tiempo o por cumplimiento, se utiliza \textit{discharge}.
\end{itemize}
\end{definition}

\begin{table}[H]
\centering
\caption{Formas de extinción de contratos en el derecho argentino}
\begin{tabularx}{\textwidth}{
    >{\raggedright\arraybackslash}p{0.28\textwidth}
    >{\raggedright\arraybackslash}X
}
\toprule
\textbf{Tipo} & \textbf{Formas} \\
\midrule
Extinción normal & Cumplimiento; vencimiento del plazo \\
\addlinespace
Extinción anormal (equivalente a \textit{terminate}) & Rescisión; resolución; revocación; nulidad; novación \\
\bottomrule
\end{tabularx}
\end{table}

\begin{note}
El término \textit{termination} puede tener distintos significados según
el contexto. En el ámbito laboral, \textit{employee termination} significa
despido de un empleado. La expresión ``\textit{the employee was terminated}''
equivale a ``el empleado fue despedido''. Es importante evaluar el contexto
para determinar el significado preciso.
\end{note}

% ============================================================
%  CAPÍTULO 4 — DISPOSICIONES OPERATIVAS
% ============================================================
\chapter{Disposiciones Operativas}

\section{Assignment and Novation (Cesión e Innovación)}

\begin{definition}{Assignment (Cesión)}
El \textbf{assignment} consiste en la cesión de los derechos y
obligaciones emergentes de un contrato a un tercero. En inglés, la
familia léxica del término incluye: \textit{assignment} (sustantivo),
\textit{assigned} (adjetivo participio) y \textit{to assign} (verbo).
\end{definition}

\begin{definition}{Novation (Innovación)}
La \textbf{novation} es la creación de una nueva obligación que se
desprende de una obligación anterior, sin dejar sin efecto la precedente.
Es un mecanismo de modificación del contrato diferente a la cesión. En
inglés: \textit{novation} (sustantivo) y \textit{novated} (adjetivo).
\end{definition}

Ambas figuras constituyen mecanismos de modificación subjetiva u objetiva
del contrato, pero operan de manera diferente. En los contratos locales
argentinos no suele incluirse esta cláusula como disposición estándar,
aunque puede aparecer según el tipo contractual.

\section{Counterparts (Ejemplares)}

\begin{definition}{Counterparts}
La cláusula de \textbf{counterparts} regula la celebración del contrato
mediante múltiples ejemplares firmados por las partes, posiblemente en
lugares distintos. Todos los ejemplares se consideran un único instrumento
con idéntico contenido (\emph{un mismo tenor}).
\end{definition}

Sus características principales son:
\begin{itemize}
    \item Se celebran dos o más ejemplares del mismo contrato.
    \item Los ejemplares pueden ser firmados a distancia por las partes.
    \item Todos los ejemplares conforman un único instrumento.
    \item El contenido es exactamente igual en todos.
    \item La finalidad es que cada parte conserve un ejemplar del acuerdo.
\end{itemize}

Esta cláusula es especialmente relevante en contratos internacionales y en
contratos de software. En el derecho argentino no se la incluye como
cláusula formal independiente; en cambio, se la incorpora en la redacción
del contrato mediante la expresión clásica: \textit{``dos ejemplares de un
mismo tenor y a un solo efecto''}.

\begin{example}{Contratos de software y términos de servicio}
Los contratos de términos de servicio que los usuarios aceptan en plataformas
digitales frecuentemente contienen, hacia el final del documento, una cláusula
de competencia que somete los conflictos a los tribunales de una jurisdicción
extranjera (por ejemplo, Nueva York). Suelen pasar inadvertidas por efecto del
asentimiento sin lectura.
\end{example}

\section{Severance Clause (Nulidad Parcial)}

\begin{definition}{Severance Clause}
La \textbf{severance clause} (cláusula de separabilidad o nulidad parcial)
establece que si una determinada cláusula del contrato es declarada nula
o inexigible (\textit{void} o \textit{impossible to enforce}), dicha
cláusula será removida del contrato \emph{sin afectar la validez del
resto del instrumento}. El verbo inglés \textit{to sever} significa
``cortar'' o ``separar''.
\end{definition}

El principio subyacente es el de \textbf{nulidad parcial}: la invalidez de
una parte del contrato no acarrea necesariamente la invalidez del todo.

\begin{important}
Esta cláusula tiene un límite: si la cláusula declarada nula es de carácter
sustancial para la economía del contrato, es posible que la nulidad afecte
al contrato en su totalidad. La separabilidad opera siempre que sea posible
y jurídicamente admisible.
\end{important}

En el sistema argentino no se incluye formalmente esta cláusula porque el
Código Civil y Comercial ya contempla el principio de nulidad parcial. No
obstante, el principio existe y puede derivarse de las normas generales del
ordenamiento.

% ============================================================
%  CAPÍTULO 5 — DISPOSICIONES ADMINISTRATIVAS
% ============================================================
\chapter{Disposiciones Administrativas}

\section{Relationship of the Parties (Naturaleza del Vínculo)}

La cláusula \textbf{relationship of the parties} tiene por objeto definir
y aclarar la naturaleza del vínculo jurídico entre las partes contratantes.
Su función principal es especificar que no existe relación de dependencia
laboral entre ellas.

Aparece con frecuencia en contratos:
\begin{itemize}
    \item Comerciales (para excluir la relación de dependencia laboral).
    \item De locación de servicios.
    \item De distribución.
    \item De prestación de servicios en general.
\end{itemize}

\begin{example}{Cláusula de naturaleza del vínculo en contratos de servicios}
Una institución puede suscribir con un docente un contrato de locación de
servicios que incluya una cláusula declarando que no existe relación de
dependencia. Sin embargo, la mera inserción de esa cláusula no determina
automáticamente la naturaleza del vínculo: si de los hechos (asistencia
semanal, horario fijo, continuidad, mismo grupo) surge una relación de
dependencia, la calificación jurídica real prevalecerá sobre la
designación formal del contrato.
\end{example}

En el derecho argentino no existe una cláusula específica con este nombre.
El vínculo entre las partes se infiere de la naturaleza del contrato y de
los hechos, con independencia de la denominación que las partes le
atribuyan.

\section{Amendment (Modificación)}

\begin{definition}{Amendment}
La cláusula de \textbf{amendment} (modificación) establece que el
contrato solo podrá ser modificado mediante acuerdo escrito de ambas
partes. Quedan prohibidas las modificaciones unilaterales: lo que fue
pactado por escrito solo puede alterarse por escrito y con consentimiento
bilateral.
\end{definition}

En el sistema argentino este principio está implícito en el principio de
\textbf{autonomía de la voluntad}, por lo que no es necesario incluirlo
como cláusula expresa. En los contratos anglosajones, en cambio, se lo
explicita por razones de seguridad jurídica.

\section{Notices (Notificaciones)}

\begin{definition}{Notices}
La cláusula de \textbf{notices} (notificaciones) regula el procedimiento
y los medios a través de los cuales las partes deben cursarse comunicaciones
formales con efectos jurídicos. Incluye habitualmente: dirección y datos
particulares de cada parte; medios admitidos (correo electrónico, fax,
servicio de mensajería); y el momento a partir del cual la notificación
se tiene por efectuada.
\end{definition}

\begin{example}{Exhaustividad de la cláusula de notificaciones en inglés}
Un contrato en inglés puede establecer que, si la notificación se entrega
por servicio nocturno (\textit{overnight}), se tendrá por efectuada en la
fecha de entrega; si se remite por correo certificado, en la fecha de
recepción; si se envía por correo electrónico, conforme a las condiciones
que el propio contrato detalle.
\end{example}

En el sistema argentino no se incluye esta cláusula como disposición
independiente. En cambio, al inicio del contrato, en la presentación de
las partes, se consignan los domicilios (real, constituido o procesal) con
la indicación de que en ellos se tendrán por válidas todas las
notificaciones. Las formas de notificación están reguladas por el propio
ordenamiento.

% ============================================================
%  CAPÍTULO 6 — ANÁLISIS LINGÜÍSTICO Y TERMINOLOGÍA
% ============================================================
\chapter{Análisis Lingüístico y Terminología Jurídica en Inglés}

\section{Colocaciones Léxicas en el Ciclo de Vida del Contrato}

Las \textbf{colocaciones léxicas} (\textit{collocations}) son combinaciones
habituales de palabras propias de un campo de especialidad. En el lenguaje
jurídico en inglés, cada etapa de la vida del contrato tiene asociadas
expresiones verbales específicas que no son intercambiables entre sí.

\begin{table}[H]
\centering
\caption{Verbos clave en el ciclo de vida de un contrato}
\begin{tabularx}{\textwidth}{
    >{\raggedright\arraybackslash}p{0.30\textwidth}
    >{\raggedright\arraybackslash}p{0.28\textwidth}
    >{\raggedright\arraybackslash}X
}
\toprule
\textbf{Etapa} & \textbf{Expresión en inglés} & \textbf{Equivalente en español} \\
\midrule
\textbf{Redacción} & Draft a contract & Redactar un contrato \\
\addlinespace
\textbf{Celebración} & Enter into a contract & Celebrar / formalizar un contrato \\
 & Sign a contract & Firmar un contrato \\
 & Execute an agreement & Celebrarlo / ejecutarlo \\
\addlinespace
\textbf{Cumplimiento} & Perform a contract & Cumplir / llevar a cabo \\
 & Comply with & Cumplir con \\
 & Apply a contract & Aplicarlo \\
\addlinespace
\textbf{Exigencia de cumplimiento} & Enforce & Hacer cumplir / exigir el cumplimiento \\
 & Force compliance & Obligar al cumplimiento \\
\bottomrule
\end{tabularx}
\end{table}

\begin{important}
El verbo \textbf{\textit{execute}} en inglés jurídico se refiere
técnicamente al momento de la \emph{celebración} del contrato, no al
de su ejecución o cumplimiento. Sin embargo, en la práctica aparece
utilizado con ambos significados. Al encontrarlo en un texto, debe
evaluarse en qué etapa de la vida del contrato se inserta para
determinar su sentido preciso.
\end{important}

\begin{example}{Interpretación de \textit{enter into a contract}}
La expresión \textit{to enter into a contract} no alude a un ingreso
físico. Se refiere al momento en que las partes formalizan y celebran el
acuerdo. Es una colocación léxica propia de la etapa de celebración
contractual.
\end{example}

\section{Verbos Modales en el Lenguaje Jurídico en Inglés}

Los \textbf{verbos modales} (auxiliares modales) son una categoría
gramatical fundamental en el lenguaje jurídico en inglés. Sus propiedades
formales son:

\begin{itemize}
    \item Son verbos auxiliares sin significado léxico propio.
    \item Necesitan ir seguidos de un verbo principal en infinitivo.
    \item Van sin la partícula \textit{to} (a diferencia de otros verbos).
    \item No se conjugan: permanecen invariables.
    \item No varían en tercera persona del singular (no reciben el morfema
    \textit{-s}).
    \item En el lenguaje jurídico expresan: obligación, potestad, permiso,
    consejo o posibilidad.
\end{itemize}

\begin{table}[H]
\centering
\caption{Verbos modales en el lenguaje jurídico en inglés}
\begin{tabularx}{\textwidth}{
    >{\raggedright\arraybackslash}p{0.14\textwidth}
    >{\raggedright\arraybackslash}p{0.30\textwidth}
    >{\raggedright\arraybackslash}X
}
\toprule
\textbf{Modal} & \textbf{Significado jurídico} & \textbf{Ejemplo} \\
\midrule
\textbf{must} & Obligación ineludible. La más fuerte. Emana de una ley o contrato. & \textit{The parties must comply} = las partes tienen la obligación de cumplir. \\
\addlinespace
\textbf{shall} & Obligación proveniente de autoridad superior. Muy utilizado en documentos legales y legislación. Poco frecuente en el lenguaje cotidiano. & \textit{The parties shall meet} = las partes deben reunirse. \\
\addlinespace
\textbf{may} & Potestad o facultad. No implica obligación, sino capacidad o derecho. & \textit{The court may decide} = el tribunal tiene la facultad de decidir. \\
\addlinespace
\textbf{might} & Posibilidad más remota. Probabilidad baja. & \textit{The claim might succeed} = el reclamo podría prosperar (con menor probabilidad). \\
\addlinespace
\textbf{should} & Recomendación o consejo. Menos imperativo que \textit{must}. & \textit{You should consult an attorney} = debería consultar a un abogado. \\
\bottomrule
\end{tabularx}
\end{table}

\begin{note}
La expresión \textbf{\textit{have to}} suele considerarse funcionalmente
equivalente a un verbo modal porque expresa obligación de manera similar
a \textit{must}. Sin embargo, técnicamente no pertenece a la categoría
gramatical de los modales puros. En contextos jurídicos, \textit{must} y
\textit{shall} son los instrumentos de obligación por excelencia.
\end{note}

\section{El Conector \textit{Despite} en Textos Jurídicos}

El conector \textbf{\textit{despite}} (a pesar de, no obstante) es un
elemento frecuente en los textos jurídicos en inglés. Introduce una
proposición subordinada que expresa concesión o contradicción con la
proposición principal. Su empleo en cláusulas contractuales introduce
excepciones, limitaciones o salvedades respecto de lo establecido en el
resto del contrato.

% ============================================================
%  CAPÍTULO 7 — IMPACTO DE LA GLOBALIZACIÓN EN LOS CONTRATOS
% ============================================================
\chapter{Impacto de la Globalización en la Contratación Argentina}

\section{La Influencia Anglosajona en los Contratos Locales}

La globalización y el crecimiento del comercio internacional han generado
una creciente incorporación de cláusulas y estructuras propias del sistema
contractual anglosajón en los contratos argentinos. Este fenómeno se
manifiesta en varios niveles:

\begin{itemize}
    \item Incorporación de terminología en inglés (\textit{junior},
    \textit{senior}, entre otros).
    \item Adopción de la estructura de cláusulas típicamente anglosajonas.
    \item Reproducción de modelos sin analizar si el Código Civil y Comercial
    ya contempla las disposiciones incluidas.
    \item Creciente presencia de la influencia del inglés en el lenguaje
    jurídico cotidiano, en documentos y en sentencias.
\end{itemize}

\begin{important}
Los contratos que resultan de la mezcla irresponsable de sistemas jurídicos
diferentes han sido denominados informalmente \textbf{``contratos
Frankenstein''}: incorporan elementos del derecho anglosajón sin considerar
que el sistema argentino ya prevé explícitamente, a través del Código Civil
y Comercial, muchas de las cláusulas que en los contratos en inglés deben
ser explicitadas. El resultado son redundancias, inconsistencias e
incompatibilidades.
\end{important}

\section{Diferencias Estructurales entre Sistemas}

Una diferencia fundamental entre el sistema jurídico anglosajón (\textit{common
law}) y el sistema argentino (de tradición romano-germánica) incide
directamente en la necesidad de explicitación de cláusulas:

\begin{itemize}
    \item En el sistema anglosajón, muchos principios no están
    necesariamente codificados. Aquello que no consta expresamente en el
    contrato puede no estar protegido. De ahí la necesidad de explicitar
    cláusulas como \textit{entire agreement}, \textit{force majeure},
    \textit{severance}, \textit{amendment} y \textit{notices}.
    \item En el sistema argentino, el Código Civil y Comercial establece
    con precisión los principios generales de los contratos: autonomía de
    la voluntad, fuerza mayor, nulidad parcial, forma de las modificaciones.
    Por ello, muchas de estas cláusulas resultan implícitas y no es
    necesario consignarlas expresamente.
\end{itemize}

\section{Recomendación Profesional}

Al redactar o revisar un contrato que incorpora cláusulas de origen
anglosajón, es necesario:
\begin{enumerate}
    \item Verificar si la disposición ya está contemplada por el
    ordenamiento argentino.
    \item Adaptar conscientemente la cláusula al sistema local, sin
    simplemente copiar y pegar modelos en inglés.
    \item Mantener la coherencia interna del contrato y la compatibilidad
    con el marco normativo aplicable.
    \item Preservar las protecciones legales sin introducir redundancias o
    incompatibilidades.
\end{enumerate}

La competencia bilingüe y el conocimiento del derecho comparado no son un
lujo profesional, sino una necesidad en un entorno de práctica jurídica
cada vez más globalizado.

% ============================================================
%  CAPÍTULO 8 — CUADRO CORNELL
% ============================================================
\chapter{Síntesis y Repaso: Método Cornell}

\begin{table}[H]
\centering
\caption{Cuadro Cornell — Cláusulas Estándar en Contratos}
\begin{tabularx}{\textwidth}{
    >{\raggedright\arraybackslash}p{0.30\textwidth}
    >{\raggedright\arraybackslash}X
}
\toprule
\textbf{Preguntas / Palabras clave} & \textbf{Notas / Respuestas} \\
\midrule

\textbf{¿Qué son las boilerplate clauses?} &
Son cláusulas contractuales estandarizadas que se replican con escasas
variaciones de un contrato a otro. Abordan cuestiones vinculadas a posibles
conflictos y suelen ubicarse al final del contrato. Frecuentemente son
pasadas por alto por las partes durante la negociación. \\
\addlinespace

\textbf{¿Cuál es el origen del término ``boilerplate''?} &
Proviene de la imprenta del siglo XIX. Las placas metálicas
(\textit{boilerplates}) contenían textos preestablecidos que se replicaban
idénticamente en cada impresión. De allí deriva la idea de texto fijo,
estandarizado y no sujeto a negociación. No debe traducirse literalmente. \\
\addlinespace

\textbf{¿Qué diferencia hay entre \textit{governing law} y \textit{jurisdiction}?} &
\textit{Governing law} (derecho aplicable) determina con qué legislación se
interpreta el contrato. \textit{Jurisdiction} en inglés suele referirse a la
competencia (qué juez entiende en la causa), no a la jurisdicción abstracta.
Ambas cláusulas deben ser coherentes entre sí. \\
\addlinespace

\textbf{¿Qué distingue jurisdicción de competencia?} &
Jurisdicción es la potestad abstracta de los jueces (todos la tienen).
Competencia determina qué juez específico debe entender en un caso,
clasificada por materia, territorio, grado y persona. La expresión
sintetizadora es: ``la competencia es la medida de la jurisdicción''. \\
\addlinespace

\textbf{¿Qué establece la cláusula \textit{entire agreement}?} &
Que solo lo consignado por escrito en el contrato es vinculante. Las
negociaciones previas y las comunicaciones anteriores a la firma no
integran el acuerdo. En Argentina es implícita por el principio de
autonomía de la voluntad, pero se incorpora cada vez más por influencia
anglosajona. \\
\addlinespace

\textbf{¿Qué diferencia hay entre \textit{force majeure} y \textit{acts of God}?} &
\textit{Acts of God} abarca exclusivamente eventos naturales sin intervención
humana. \textit{Force majeure} es el concepto más amplio: incluye tanto eventos
naturales como de origen humano (guerras, huelgas, disturbios). Ambos
equivalen en términos generales a la fuerza mayor y el caso fortuito del
derecho argentino. \\
\addlinespace

\textbf{¿Qué diferencia hay entre \textit{terminate} y \textit{discharge}?} &
\textit{Terminate} implica poner fin anticipado al contrato (extinción anormal:
rescisión, resolución, revocación). \textit{Discharge} es el concepto más
general y abarca también la extinción normal por cumplimiento o vencimiento
del plazo. En el ámbito laboral, \textit{termination} equivale a despido. \\
\addlinespace

\textbf{¿Qué es la \textit{severance clause}?} &
Es la cláusula de nulidad parcial. Establece que si una cláusula es declarada
nula o inexigible, será removida del contrato sin afectar la validez del
resto. Tiene como límite el caso en que la cláusula afectada sea sustancial
para el contrato, en cuyo caso puede comprometer la validez total. \\
\addlinespace

\textbf{¿Qué son los \textit{counterparts}?} &
Múltiples ejemplares de un mismo contrato firmados a distancia por las partes.
Todos se consideran un único instrumento con idéntico contenido. En Argentina
se incorpora en la redacción mediante la fórmula ``dos ejemplares de un mismo
tenor y a un solo efecto''. \\
\addlinespace

\textbf{¿Qué diferencia hay entre \textit{assignment} y \textit{novation}?} &
\textit{Assignment} es la cesión de derechos y obligaciones emergentes del
contrato a un tercero. \textit{Novation} es la creación de una nueva
obligación que se desprende de la anterior, sin necesariamente dejar sin
efecto la original. Son dos mecanismos distintos de modificación contractual. \\
\addlinespace

\textbf{¿Para qué sirve la cláusula \textit{relationship of the parties}?} &
Define y aclara la naturaleza del vínculo jurídico entre las partes,
especificando en particular que no existe relación de dependencia laboral.
Es frecuente en contratos de prestación de servicios y locación. En
Argentina no existe como cláusula formal independiente. \\
\addlinespace

\textbf{¿Cuál es la función de la cláusula \textit{notices}?} &
Regula el procedimiento y los medios para las comunicaciones formales con
efectos jurídicos: dirección, medios admitidos (email, fax, mensajería) y
el momento a partir del cual la notificación se considera efectuada. En
Argentina se incorpora en el encabezado del contrato al presentar a las partes. \\
\addlinespace

\textbf{¿Cómo se usa \textit{must} vs. \textit{shall} en lenguaje jurídico?} &
\textit{Must} indica obligación ineludible emanada de una ley o contrato.
\textit{Shall} también expresa obligación, con connotación de autoridad
superior; es el modal jurídico por excelencia, muy utilizado en documentos
legales y legislación pero poco frecuente en el lenguaje cotidiano. \\
\addlinespace

\textbf{¿Cómo se distinguen \textit{may} y \textit{might}?} &
\textit{May} expresa potestad o facultad (probabilidad alta o capacidad legal).
\textit{Might} expresa posibilidad más remota (probabilidad menor). En
documentos jurídicos, \textit{may} suele conferir una facultad discrecional
al sujeto al que se refiere. \\
\addlinespace

\textbf{¿Qué son los ``contratos Frankenstein''?} &
Contratos que resultan de incorporar cláusulas del sistema anglosajón en
contratos argentinos sin verificar si el Código Civil y Comercial ya las
contempla. Generan redundancias, inconsistencias e incompatibilidades entre
sistemas jurídicos diferentes. \\

\midrule

\multicolumn{2}{p{\textwidth}}{
\textbf{Resumen:}
Las cláusulas estándar en contratos (\textit{boilerplate clauses}) son
disposiciones estandarizadas cuya denominación genérica suele llevar a que
sean pasadas por alto. Regulan aspectos clave como el derecho aplicable y la
competencia (que deben estar alineados), la integración del contrato
(\textit{entire agreement}), la fuerza mayor (\textit{force majeure}) y los
\textit{acts of God}, las formas de extinción anticipada (\textit{terminate}
vs. \textit{discharge}), la cesión e innovación, la nulidad parcial
(\textit{severance}), las notificaciones y la naturaleza del vínculo. En el
plano lingüístico, las colocaciones léxicas asociadas a cada etapa del
contrato y los verbos modales (\textit{must}, \textit{shall}, \textit{may},
\textit{might}, \textit{should}) son herramientas de precisión jurídica. La
globalización genera una creciente incorporación de estas cláusulas en
contratos argentinos, lo que exige una adaptación consciente al sistema
local para evitar los denominados ``contratos Frankenstein''.
} \\

\bottomrule
\end{tabularx}
\end{table}

\end{document}
